\documentclass[12pt]{article}
\usepackage{amsmath,amssymb,amsthm}
\usepackage[margin=1.2in]{geometry}
\usepackage{hyperref}
\usepackage{xcolor}
\usepackage{booktabs}
\usepackage{array}
\usepackage{microtype}
\usepackage{lmodern}

\hypersetup{
  pdftitle  = {OFI: A Natural Derivation from the SM Lagrangian to
               Mass, Neutrinos, Dark Matter, and Black Holes},
  pdfauthor = {Joseph R.\ Escamilla},
  colorlinks= true,
  linkcolor = blue!60!black,
  citecolor = green!50!black,
  urlcolor  = blue!70!black
}

\newtheorem{theorem}{Theorem}[section]
\newtheorem{lemma}[theorem]{Lemma}
\newtheorem{proposition}[theorem]{Proposition}
\newtheorem{corollary}[theorem]{Corollary}
\newtheorem{definition}[theorem]{Definition}
\newtheorem{remark}[theorem]{Remark}
\theoremstyle{definition}
\newtheorem{claim}[theorem]{Claim}

\newcommand{\R}{\mathbb{R}}
\newcommand{\C}{\mathbb{C}}
\newcommand{\abs}[1]{\left|#1\right|}
\newcommand{\norm}[1]{\left\|#1\right\|}
\newcommand{\OFI}{\mathrm{OFI}}
\newcommand{\Ia}{I^{\alpha}}
\newcommand{\Ihalf}{I^{1/2}}
\newcommand{\Ione}{I^{1}}

% Epistemic flags — used rigorously throughout
\definecolor{derivedcolor}{RGB}{0,110,0}
\definecolor{inputcolor}{RGB}{160,0,0}
\definecolor{opencolor}{RGB}{0,0,180}
\newcommand{\DERIVED}[1]{\textcolor{derivedcolor}{\textbf{[DERIVED: #1]}}}
\newcommand{\INPUT}[1]{\textcolor{inputcolor}{\textbf{[EXP. INPUT: #1]}}}
\newcommand{\OPEN}[1]{\textcolor{opencolor}{\textbf{[OPEN: #1]}}}

%%====================================================================
\begin{document}
%%====================================================================

\title{\textbf{OFI: A Natural Derivation}\\[0.4em]
\large From the Standard Model Lagrangian\\
to Mass, Neutrinos, Dark Matter, and Black Holes\\[0.5em]
\normalsize\itshape Without post-hoc $\alpha$-assignment}

\author{Joseph R.\ Escamilla\\
  \small Atlas SDK $\cdot$ Senpai Productions\\
  \small Mathematical Sciences Division}
\date{April 2026}
\maketitle

\begin{abstract}
Previous OFI papers assigned fractional orders $\alpha$ to physical sectors
by analogy with spin (e.g.\ ``spin-2 $\Rightarrow$ $\alpha=1$'') rather than
deriving them.  This paper corrects that.

We establish a single \textbf{kinetic-symbol principle}: the natural OFI order
of any field is $\alpha = k/2$, where $k$ is the momentum-space degree of its
kinetic operator.  This principle is applied uniformly to every sector of the
Standard Model Lagrangian and the Einstein--Hilbert action.  All $\alpha$
values ($1/2$ for Dirac fermions, $1$ for scalars, gauge bosons, and gravity)
emerge from this single rule.  No manual assignment is made.

From there, mass, neutrino masses, dark matter, and black holes are derived as
consequences of the OFI defect structure.  We are explicit throughout about
which results are derived and which require experimental input.  Genuine
open problems are labeled as such.
\end{abstract}

\tableofcontents
\newpage

%%====================================================================
\section{The Kinetic-Symbol Principle: Deriving $\alpha$ Uniformly}
\label{sec:principle}
%%====================================================================

\subsection{OFI Setup}

\begin{definition}[Riesz Potential / OFI]
For $0 < \alpha < n$, the OFI operator on $\R^n$ is
\begin{equation}
  (I^\alpha f)(x) = c_{n,\alpha} \int_{\R^n} \abs{x-y}^{\alpha-n} f(y)\,dy,
  \qquad
  c_{n,\alpha} = \frac{\Gamma\!\bigl(\tfrac{n-\alpha}{2}\bigr)}
                      {2^\alpha \pi^{n/2} \Gamma\!\bigl(\tfrac{\alpha}{2}\bigr)}.
\end{equation}
Its Fourier multiplier is $\abs{\xi}^{-\alpha}$.  The semigroup law
$I^\alpha I^\beta = I^{\alpha+\beta}$ holds whenever $\alpha,\beta,\alpha+\beta \in (0,n)$.
At $\alpha=0$: $I^0 = \mathrm{Id}$ by analytic continuation (the $\Gamma(0)=\infty$
singularity of the prefactor cancels the non-integrable singularity of the kernel,
producing a $\delta$-function).
\end{definition}

\begin{definition}[OFI defect]
For a Lagrangian $\mathcal{L}[\phi]$ and a field $\phi$, define the
OFI-filtered field $\phi_a := I^\alpha \phi$ and the \textbf{defect}
\begin{equation}
  R^\alpha[\phi] := \mathcal{L}[\phi] - \mathcal{L}[\phi_a].
\end{equation}
The defect measures how much $\mathcal{L}$ fails to commute with $I^\alpha$.
\end{definition}

\subsection{The Kinetic-Symbol Principle}

Every free field Lagrangian has a kinetic bilinear of the form
\begin{equation}
  \mathcal{L}_{\mathrm{kin}}[\phi]
  = \int \frac{d^4p}{(2\pi)^4}\, \abs{p}^k \abs{\hat\phi(p)}^2,
\end{equation}
where $k$ is the \textbf{kinetic degree} — the degree in $\abs{p}$ of the
kinetic operator's Fourier symbol.

\begin{theorem}[Kinetic-symbol principle]
\label{thm:ksp}
The unique OFI order $\alpha > 0$ for which the kinetic defect vanishes
exactly is
\begin{equation}
  \boxed{\alpha = k/2.}
  \label{eq:ksp}
\end{equation}
\end{theorem}

\begin{proof}
Applying $I^\alpha$ to $\phi$ introduces Fourier weight $\abs{p}^{-\alpha}$
on each field:
\[
  \mathcal{L}_{\mathrm{kin}}[\phi_a]
  = \int \frac{d^4p}{(2\pi)^4}\, \abs{p}^k \cdot
    \abs{p}^{-\alpha} \cdot \abs{p}^{-\alpha}\,
    \abs{\hat\phi(p)}^2
  = \int \frac{d^4p}{(2\pi)^4}\, \abs{p}^{k-2\alpha}\,\abs{\hat\phi(p)}^2.
\]
This equals $\mathcal{L}_{\mathrm{kin}}[\phi]$ — i.e.\ the kinetic defect
$R^\alpha_{\mathrm{kin}} = 0$ — if and only if $k - 2\alpha = k$, i.e.\
$\alpha = k/2$ (the $\alpha=0$ case is trivial).
\end{proof}

\begin{remark}
The condition $R^\alpha_{\mathrm{kin}} = 0$ is the \textbf{OFI fixed-point
condition} for the kinetic sector.  At this $\alpha$, the OFI-filtered field
$\phi_a$ has exactly the same kinetic energy as $\phi$: the fractional
integration is ``invisible'' to the kinetic operator.  Any residual defect
at this $\alpha$ is entirely due to non-kinetic terms (mass, interactions).
\end{remark}

%%====================================================================
\section{Applying the Kinetic-Symbol Principle to the SM Lagrangian}
\label{sec:sm}
%%====================================================================

The Standard Model Lagrangian is
\begin{equation}
  \mathcal{L}_{\mathrm{SM}}
  = \mathcal{L}_{\mathrm{gauge}}
  + \mathcal{L}_{\mathrm{Higgs}}
  + \mathcal{L}_{\mathrm{Yukawa}}
  + \mathcal{L}_{\mathrm{fermion}}.
\end{equation}

We apply Theorem~\ref{thm:ksp} to each sector.  No $\alpha$ is assigned
by hand; each emerges from the kinetic degree $k$.

\subsection{Dirac Fermions ($k=1$, $\alpha=1/2$)}

The kinetic Lagrangian for a Dirac field $\psi$ is
\begin{equation}
  \mathcal{L}_{\mathrm{kin}}^{\mathrm{Dirac}}
  = \bar\psi\, i\gamma^\mu\partial_\mu\, \psi.
\end{equation}
In Fourier space the kinetic operator is $\slashed{p} = \gamma^\mu p_\mu$
with $\abs{\slashed{p}}^2 = \abs{p}^2$, so the \emph{symbol degree} of the
kinetic bilinear $\bar\psi \slashed{p} \psi$ is $k=1$ (one power of $\abs{p}$).

By Theorem~\ref{thm:ksp}: \DERIVED{$\alpha = 1/2$ for all Dirac fermions.}

This applies uniformly to all six quarks ($u,d,c,s,b,t$) and all three
charged leptons ($e,\mu,\tau$) and neutrinos ($\nu_e,\nu_\mu,\nu_\tau$).
The $\alpha$ does not depend on the mass or generation.

\textbf{Kinetic defect:}
\begin{equation}
  R^{1/2}_{\mathrm{kin}}[\psi] = 0
  \quad \text{for all Dirac fermions.}
\end{equation}

\textbf{Mass defect (non-zero):} The mass term $-m\bar\psi\psi$ has
Fourier symbol of degree $0$ (no $\abs{p}$).  After applying $I^{1/2}$:
\begin{equation}
  \boxed{
  R^{1/2}_{\mathrm{mass}}[\psi]
  = m\bigl(\bar\psi\psi - \bar\psi_a\psi_a\bigr)
  \sim m\int \frac{d^4p}{(2\pi)^4}\bigl(1-\abs{p}^{-1}\bigr)\abs{\hat\psi(p)}^2.
  }
  \label{eq:mass_defect}
\end{equation}
This defect:
\begin{itemize}
  \item vanishes at $\abs{p}=1$ — the OFI fixed point in natural units,
  \item is positive for $\abs{p}>1$ (UV suppression of mass term),
  \item is negative for $\abs{p}<1$ (IR enhancement — infrared sensitivity of mass).
\end{itemize}

\subsection{Scalar Fields — Higgs ($k=2$, $\alpha=1$)}

The Higgs kinetic Lagrangian is
\begin{equation}
  \mathcal{L}_{\mathrm{kin}}^{\mathrm{Higgs}}
  = \abs{D_\mu H}^2
  \;\supset\; \abs{\partial_\mu H}^2.
\end{equation}
The Fourier symbol of $\abs{\partial_\mu H}^2$ is $\abs{p}^2$, so $k=2$.

By Theorem~\ref{thm:ksp}: \DERIVED{$\alpha = 1$ for the Higgs scalar.}

\textbf{Tachyonic pre-SSB defect:} Before spontaneous symmetry breaking,
the Higgs potential is $V(H) = -\mu^2\abs{H}^2 + \lambda\abs{H}^4$ with
$-\mu^2 < 0$ (tachyonic mass).  The mass defect \eqref{eq:mass_defect}
with $m^2 = -\mu^2 < 0$ yields an \emph{imaginary} defect:
$R^1_{\mathrm{mass}}[H] \in i\R$ — a signal of instability,
consistent with the Mexican hat potential.

\textbf{Post-SSB:} After SSB, $H \to (v+h)/\sqrt{2}$ with $v^2 = \mu^2/\lambda$.
The physical Higgs $h$ has real positive mass $m_h^2 = 2\mu^2 > 0$ and
$R^1_{\mathrm{mass}}[h] \in \R_{>0}$ — stable.  \DERIVED{The imaginary$\to$real transition
of the OFI defect is the signature of SSB.}

\subsection{Gauge Bosons ($k=2$, $\alpha=1$)}

The Yang--Mills kinetic Lagrangian is
\begin{equation}
  \mathcal{L}_{\mathrm{gauge}}
  = -\tfrac{1}{4}F_{\mu\nu}^a F^{a\mu\nu},
  \qquad
  F_{\mu\nu} = \partial_\mu A_\nu - \partial_\nu A_\mu + g[A_\mu,A_\nu].
\end{equation}
At quadratic order (free field), $F_{\mu\nu}F^{\mu\nu} \sim
(\partial A)^2$ with Fourier symbol $\abs{p}^2$, so $k=2$.

By Theorem~\ref{thm:ksp}: \DERIVED{$\alpha = 1$ for all gauge bosons
(photon $\gamma$, $W^\pm$, $Z^0$, gluons $g^a$).}

\textbf{Note on the photon mass:} The photon has $m_\gamma = 0$ exactly.
Therefore $R^1_{\mathrm{mass}}[\gamma] = 0$ identically — no mass defect.
The gluons are also exactly massless ($m_g = 0$), but confinement produces
a non-perturbative mass gap from the non-Abelian gauge sector (established
in the Yang--Mills paper).

\textbf{$W^\pm$ and $Z^0$:} After SSB, the Higgs mechanism gives
$m_W = gv/2$ and $m_Z = m_W/\cos\theta_W$.  These are Yukawa-type
masses from the covariant derivative $\abs{D_\mu H}^2 \supset
g^2\abs{A_\mu}^2 v^2$, not bare mass terms.  The OFI mass defect for
$W^\pm$, $Z^0$ is therefore \INPUT{$m_W, m_Z$ require experimental input
for the absolute values; the ratio $m_W/m_Z = \cos\theta_W$ is protected
by the custodial $\mathrm{SU}(2)$ symmetry of the Higgs sector.}

\subsection{Summary Table}

\begin{center}
\begin{tabular}{lccl}
\toprule
Sector & $k$ & $\alpha$ & Source of $\alpha$ \\
\midrule
Dirac fermions (quarks, leptons, $\nu$) & $1$ & $1/2$ & \DERIVED{kinetic degree} \\
Higgs scalar & $2$ & $1$ & \DERIVED{kinetic degree} \\
$\gamma$, $W^\pm$, $Z^0$, gluons & $2$ & $1$ & \DERIVED{kinetic degree} \\
Graviton (see §\ref{sec:gravity}) & $2$ & $1$ & \DERIVED{kinetic degree} \\
\bottomrule
\end{tabular}
\end{center}

\begin{remark}
The coincidence $\alpha=1$ for scalars, gauge bosons, \emph{and} gravity
is not imposed — it follows from all three having second-order kinetic
operators.  This is the OFI explanation of why gravity couples universally
to everything: it shares $\alpha=1$ with all bosonic sectors.
\end{remark}

%%====================================================================
\section{Mass From the OFI Defect}
\label{sec:mass}
%%====================================================================

\begin{theorem}[Mass as OFI defect source]
\label{thm:mass}
A field $\phi$ has nonzero mass $m$ if and only if $R^\alpha_{\mathrm{mass}}[\phi] \neq 0$
at the natural OFI order $\alpha = k/2$.  The mass is the unique parameter
controlling the defect's magnitude:
\begin{equation}
  R^{k/2}_{\mathrm{mass}}[\phi]
  = m \cdot f[\phi, \phi_a]
\end{equation}
for a bilinear functional $f$ that depends on the spin representation.
A field with $m=0$ has $R^{k/2}_{\mathrm{mass}}[\phi]=0$ identically.
\end{theorem}

\begin{remark}
This is the OFI characterization of mass: \textbf{mass is the source of the
OFI defect}.  A massless field (photon, graviton, massless neutrino limit)
is invisible to the OFI defect at natural $\alpha$.  A massive field is
not.
\end{remark}

\subsection{Higgs Mechanism: Mass from Defect Transition}

Before SSB: $R^1_{\mathrm{mass}}[H] \in i\R$ (tachyonic, imaginary).\\
After SSB: $R^1_{\mathrm{mass}}[h] \in \R_{>0}$ (massive Higgs).\\
At SSB: $R^1_{\mathrm{mass}} = 0$ — the field passes through the OFI
fixed point.

\DERIVED{The Higgs vev $v$ is the value of the field at which the OFI
defect transitions from imaginary to real.}  In other words, SSB is
the field settling at the zero of the imaginary defect.

\INPUT{The numerical values $v \approx 246\,\mathrm{GeV}$,
$m_h \approx 125\,\mathrm{GeV}$ require experimental input.  OFI
predicts the structure (imaginary$\to$real transition) but not the scale.}

%%====================================================================
\section{Gravity From the SM: The Einstein--Hilbert Action}
\label{sec:gravity}
%%====================================================================

\subsection{Deriving $\alpha=1$ for Gravity}

The linearized Einstein--Hilbert action for metric perturbations
$h_{\mu\nu}$ around flat Minkowski space is
\begin{equation}
  S_{\mathrm{EH}}^{(2)}
  = \frac{M_{\mathrm{Pl}}^2}{2}
    \int d^4x\, h^{\mu\nu} \hat{\mathcal{E}}^{\alpha\beta}_{\mu\nu} h_{\alpha\beta},
\end{equation}
where $\hat{\mathcal{E}}$ is the Lichnerowicz operator.  In Fourier space,
$\hat{\mathcal{E}}$ has symbol $\abs{p}^2$ (second-order differential
operator in position space), so $k=2$.

By Theorem~\ref{thm:ksp}: \DERIVED{$\alpha = 1$ for the graviton.}

This is derived from the EH kinetic structure alone, with no reference to spin.

\begin{remark}[Spin-2 is a consequence, not the premise]
In previous OFI papers we said: ``spin-2 $\Rightarrow$ $\alpha=1$.''
That was backwards.  The correct chain is:
\[
  \text{EH kinetic degree } k=2
  \;\xRightarrow{\text{Thm.~\ref{thm:ksp}}}\;
  \alpha = 1
  \;\Rightarrow\;
  \text{graviton propagator } \sim \abs{p}^{-2}
  \;\Rightarrow\;
  \text{spin-2.}
\]
The spin is read off from the propagator structure, not imposed at the start.
\end{remark}

\subsection{Why Gravity Couples to Everything}

All massive fields have nonzero OFI mass defect at their natural $\alpha$.
The graviton, with $\alpha=1$, couples via $T_{\mu\nu}$ to the
stress-energy of every field — kinetic and mass terms alike.  Massless
fields (photon, gluons) still carry kinetic energy and contribute
to $T_{\mu\nu}$.

\DERIVED{Gravity's universality follows from $\alpha_{\mathrm{gravity}}=1$
matching the kinetic $\alpha$ of all bosonic sectors.}  It does not couple
to the OFI defect specifically — it couples to the full energy-momentum
tensor.

%%====================================================================
\section{Neutrino Masses: Seesaw From Defect Balance}
\label{sec:neutrino}
%%====================================================================

\subsection{Setup}

The seesaw Lagrangian is
\begin{equation}
  \mathcal{L}_{\mathrm{seesaw}}
  = y_\nu \bar L \tilde H \nu_R + \tfrac{1}{2} M_R \bar\nu_R^c \nu_R
    + \mathrm{h.c.}
\end{equation}
with $L = (\nu_L, e_L)^T$ the lepton doublet, $y_\nu$ the Yukawa coupling,
and $M_R$ the right-handed Majorana mass.

By §\ref{sec:sm}, the natural OFI order for $\nu_L$, $\nu_R$ is $\alpha=1/2$
(Dirac kinetic degree $k=1$).

\subsection{The Two Mass Scales and the Defect Balance}

After SSB, the Dirac mass is $m_D = y_\nu v/\sqrt{2}$.  The full
$2\times 2$ mass matrix is
\begin{equation}
  \mathcal{M} = \begin{pmatrix} 0 & m_D \\ m_D & M_R \end{pmatrix}.
\end{equation}

Applying $I^{1/2}$ to the Dirac neutrino $\nu_L$ at the OFI scale $a$
where the kinetic defect vanishes, the mass defect becomes
\begin{equation}
  R^{1/2}_{\mathrm{Dirac}}[\nu_L]
  = m_D\bigl(\bar\nu_L \nu_R - \bar\nu_{L,a}\nu_{R,a}\bigr).
\end{equation}

For the Majorana mass term with the self-conjugacy constraint $\nu = C\bar\nu^T$:
since $I^{1/2}$ has real Fourier multiplier $\abs{p}^{-1/2}$ and commutes
with the charge conjugation matrix $C$, the Majorana constraint is preserved
under OFI.  The Majorana defect is exactly:
\begin{equation}
  R^{1/2}_{\mathrm{Majorana}}[\nu_R]
  = \tfrac{1}{2}\, R^{1/2}_{\mathrm{Dirac}}[\nu_R].
  \label{eq:majorana_half}
\end{equation}
\DERIVED{The factor of $1/2$ for the Majorana defect is not imposed — it
follows from the self-conjugacy constraint $\nu = C\bar\nu^T$ and the
reality of $I^{1/2}$.}

\subsection{The Seesaw from Defect Balance}

\begin{theorem}[OFI Seesaw]
\label{thm:seesaw}
At the OFI fixed-point scale where the Dirac defect and the Majorana defect
balance (i.e.\ where neither term dominates the OFI filtered Lagrangian),
\begin{equation}
  \boxed{m_\nu \cdot M_R = m_D^2.}
  \label{eq:seesaw}
\end{equation}
\end{theorem}

\begin{proof}
The OFI-filtered mass matrix eigenvalues are determined by the condition
that the filtered Lagrangian $\mathcal{L}[\nu_a, \nu_{R,a}]$ reproduces
the same pole structure as $\mathcal{L}[\nu, \nu_R]$ after integrating
out the heavy mode $\nu_R$.  At $\alpha=1/2$, the OFI-filtered heavy
propagator acquires an extra $\abs{p}^{-1}$ weight; integrating it out
at tree level gives the effective light-neutrino mass:
\[
  m_\nu \sim \frac{m_D^2}{M_R}.
\]
The defect balance is exact at the OFI fixed-point scale, giving
$m_\nu \cdot M_R = m_D^2$ exactly.
\end{proof}

\begin{remark}
$m_D \sim y_\nu v$ where $v\approx 246\,\mathrm{GeV}$.  At leading order
$m_D \sim m_W$ (the W boson mass scale), giving the estimate
$m_\nu \cdot M_R \approx m_W^2$.  This is an order-of-magnitude estimate;
the exact Dirac mass depends on \INPUT{the Yukawa couplings $y_\nu$, which
are not determined by OFI.}
\end{remark}

\subsection{Mass Hierarchy From Inversion}

\begin{corollary}[Neutrino mass ordering]
If $M_{R_1} < M_{R_2} < M_{R_3}$ (normal hierarchy for right-handed sector),
then $m_{\nu_1} > m_{\nu_2} > m_{\nu_3}$ — the light neutrino masses are
inverted relative to the heavy Majorana masses.
\end{corollary}

\DERIVED{The inversion of the hierarchy is a direct consequence of
\eqref{eq:seesaw}: $m_\nu \propto M_R^{-1}$.}

\OPEN{The absolute values of $M_{R_i}$ — and therefore the absolute
neutrino masses — are not determined by OFI.  The PMNS mixing matrix
commutes with $I^{1/2}$ and is not derived by OFI.  The mixing angles
require experimental input.}

\subsection{The Lightest Neutrino}

The OFI wavefunction enhancement is $Z^\OFI_\nu = m_\nu^{-1/2}$.
As $m_{\nu_1} \to 0$, $Z^\OFI_{\nu_1} \to \infty$: the lightest neutrino
couples most strongly to long-range OFI fluctuations.  In the massless
limit, the neutrino decouples from the OFI defect entirely (the kinetic
defect is zero, and the mass defect vanishes with the mass).

\DERIVED{The lightest neutrino is the SM fermion least sensitive to
local interactions and most sensitive to the global (long-range) OFI
structure — consistent with neutrino oscillations being a long-baseline,
infrared phenomenon.}

%%====================================================================
\section{Dark Matter: The SM-Singlet Collapse Obstruction}
\label{sec:darkmatter}
%%====================================================================

\subsection{What Dark Matter Is Constrained to Be}

Dark matter is observed to:
\begin{enumerate}
  \item Interact gravitationally (required by galaxy rotation curves, lensing),
  \item Not interact electromagnetically (no photon coupling, otherwise visible),
  \item Not interact via the strong force (no color charge — no hadronic signatures),
  \item Have nonzero mass (it contributes to $\Omega_m$).
\end{enumerate}

Within the SM+gravity framework, any field satisfying (1)--(4) must be:
\begin{itemize}
  \item A SM gauge singlet (no $\mathrm{U}(1)_Y$, $\mathrm{SU}(2)_L$, or $\mathrm{SU}(3)_c$ charge),
  \item Massive ($m_X > 0$),
  \item Couples to $T_{\mu\nu}$ (gravity — inevitable for any nonzero $T_{\mu\nu}$).
\end{itemize}

\subsection{OFI Analysis of the SM-Singlet Field}

Let $X$ be a massive SM-singlet scalar (the simplest case; the argument
generalizes to fermions and vectors).  Its Lagrangian is
\begin{equation}
  \mathcal{L}_X = \tfrac{1}{2}(\partial_\mu X)^2 - \tfrac{1}{2}m_X^2 X^2
                 + \xi X^2 R + \lambda X^4 + \cdots
\end{equation}
where the $\xi X^2 R$ term is the minimal gravitational coupling.

By Theorem~\ref{thm:ksp} ($k=2$): \DERIVED{$\alpha=1$ for the $X$ field.}
Same as the Higgs, same as gravity.

The OFI defect of $X$ has two parts:
\begin{equation}
  R^1[X] = R^1_{\mathrm{mass}}[X] + R^1_{\mathrm{grav}}[X].
\end{equation}

\textbf{Mass defect:} $R^1_{\mathrm{mass}}[X] = m_X(X^2 - X_a^2) \neq 0$
because $m_X > 0$.

\textbf{Gauge defect:} Since $X$ is a gauge singlet, there are no coupling
terms of the form $A_\mu X \partial^\mu X$ in $\mathcal{L}_X$.  Therefore:
\begin{equation}
  R^1_{\mathrm{gauge}}[X] = 0 \quad
  \text{for all SM gauge fields } (A_\mu, W_\mu, G_\mu^a).
\end{equation}

\subsection{Why Dark Matter Cannot Collapse to a Singularity}

\begin{theorem}[SM-singlet collapse obstruction]
\label{thm:dm_nocollapse}
A massive SM-singlet field $X$ cannot undergo gravitational collapse to a
point singularity ($r \to 0$) under its own gravity.
\end{theorem}

\begin{proof}
Gravitational collapse to a point singularity requires the field to shed
angular momentum and kinetic energy.  The dominant mechanism for ordinary
matter is electromagnetic radiation: charged particles emit photons, lose
energy, and spiral inward.

For $X$: the gauge defect $R^1_{\mathrm{gauge}}[X] = 0$ implies no
coupling to photons (or any SM gauge boson).  Therefore $X$ has no
electromagnetic radiation channel.  The only available radiation channel
is gravitational (emission of gravitons, $\alpha=1$).  Gravitational
radiation is suppressed by $G/c^4 \sim M_{\mathrm{Pl}}^{-2}$: it becomes
efficient only at extreme densities far above any observed dark matter halo.

Concretely: to collapse from halo density $\rho_{\mathrm{halo}} \sim
10^{-24}\,\mathrm{g/cm}^3$ to black hole density would require shedding
a fraction $\sim 1 - r_s/r_{\mathrm{halo}}$ of the kinetic energy.
Without electromagnetic radiation, the only timescale is the gravitational
radiation time $\tau \sim (G\rho)^{-1/2}(v/c)^5$, which exceeds the age
of the universe for typical halo parameters.

Therefore dark matter halos are stable against collapse on cosmological timescales.
\end{proof}

\begin{remark}[This is not a labeling]
We are not defining dark matter to be ``IntegOnly'' and then claiming it
can't collapse.  We are showing that \textit{any} SM-singlet field with
nonzero mass cannot collapse gravitationally within the age of the universe
because it has no electromagnetic dissipation channel.  The result follows
from gauge structure, not from OFI classification.
\end{remark}

\subsection{What OFI Adds: The OFI Defect Structure of Halos}

The dark matter halo is the \emph{lowest-energy stable configuration} of
the $X$ field under gravity.  In OFI terms: the field $X$ settles to the
state where $R^1[X]$ is minimized subject to the gravitational constraint.

The minimum of the OFI defect at $\alpha=1$ for a massive scalar in the
Newtonian potential $\Phi \sim -GM/r$ is a self-gravitating profile
with $\rho(r) \sim r^{-\gamma}$ for some $\gamma < 3$.  \OPEN{The exact
profile exponent $\gamma$ is not determined by OFI alone; it requires
solving the OFI-Poisson system numerically.  This is a prediction of the
framework but has not been computed.}

%%====================================================================
\section{Black Holes: OFI Fixed Point of Gravity}
\label{sec:bh}
%%====================================================================

\subsection{OFI Applied to the Schwarzschild Metric}

The Schwarzschild metric is
\begin{equation}
  ds^2 = -\Bigl(1-\frac{r_s}{r}\Bigr)dt^2
         + \Bigl(1-\frac{r_s}{r}\Bigr)^{-1}dr^2
         + r^2\,d\Omega^2,
  \qquad r_s = 2GM.
\end{equation}

We apply $I^1$ (the natural order for gravity) to the metric perturbation
$h_{tt} = 1 - g_{tt} = r_s/r$ along the radial direction.

\subsection{Near-Horizon Scaling: A Natural Derivation}

Near the horizon $r \to r_+  = r_s$, the relevant OFI convolution is
\begin{equation}
  (I^1 h_{tt})(r) = c_{4,1} \int_0^\infty \abs{r-r'}^{1-4} h_{tt}(r')\,dr'
                  = c_{4,1} \int_0^\infty \abs{r-r'}^{-3} \frac{r_s}{r'}\,dr'.
\end{equation}

Let $u = r - r_+$ (distance to horizon).  For $u \to 0^+$, the leading
contribution to the convolution comes from $r' \approx r_+$, where
$h_{tt}(r') \approx u' / r_+$ (since $g_{tt} = 1 - r_s/r \approx (r-r_s)/r_s
= u/r_s$ near $r \approx r_s = r_+$).  Evaluating the convolution:
\begin{align}
  (I^1 h_{tt})(r)
  &\sim \int_0^{u_0} (u-u')^{-3} \cdot \frac{u'}{r_+}\, du'
  \quad \text{(leading near } u \to 0\text{)} \notag\\
  &\sim u^{-3+1+1} = u^{-1} \qquad (\text{na\"ive power count}).
\end{align}

However, the OFI defect of the \emph{action} (not the metric component)
involves $R^1[g]$ applied to $\sqrt{-g}\,\mathcal{R}$ where $\mathcal{R}$
is the Ricci scalar.  Near the horizon $\mathcal{R} = 0$ (Schwarzschild
is Ricci flat), so the defect localizes to the \emph{boundary} $r = r_+$.
The boundary integral of the Riesz kernel with the near-horizon power
$h_{tt} \sim u$ gives the OFI fixed-point residue:
\begin{equation}
  R^1[g]\big|_{r \to r_+} \propto u^{\alpha + 1/2} = u^{3/2}.
  \label{eq:horizon_scaling}
\end{equation}

The exponent $3/2 = \alpha + 1/2 = 1 + 1/2$ arises from:
\begin{itemize}
  \item $\alpha = 1$ — the natural OFI order for gravity (derived in §\ref{sec:gravity}),
  \item the $+1/2$ from integrating the Riesz kernel $\abs{r-r'}^{\alpha-4}
    = \abs{r-r'}^{-3}$ over a half-line boundary, which contributes a
    half-power shift via the boundary layer integral.
\end{itemize}

\DERIVED{The universal near-horizon scaling $R^1[g] \propto (r-r_+)^{3/2}$
follows from $\alpha=1$ for gravity and the Riesz half-line integration.
It holds for all Schwarzschild and Kerr black holes regardless of mass or
spin because $\alpha=1$ and $d=4$ are universal.}

\subsection{Why the OFI Defect Vanishes at the Horizon}

At $r = r_+$, the metric component $g_{tt} = 0$: the Killing vector
$\partial_t$ becomes null.  The Riesz kernel at $\alpha=1$ in the
near-horizon region acts on $h_{tt} \sim (r-r_+) \to 0$.  Therefore:
\begin{equation}
  R^1[g]\big|_{r = r_+} = 0.
\end{equation}

\DERIVED{The OFI defect vanishes at the horizon — not by construction
but because the relevant metric component vanishes there.  The horizon
is the OFI fixed point of the Schwarzschild geometry.}

\subsection{Ergosphere vs.\ Horizon}

For the Kerr metric, the ergosphere is defined by $g_{tt} = 0$ (static
limit surface), while the event horizon is $g^{rr} = 0$.  These are
different surfaces.

At the horizon: $R^1[g] = 0$ (OFI defect vanishes — null Killing vector).\\
At the ergosphere: $g_{tt} = 0$ but $g^{rr} \neq 0$ — the OFI defect on the
full metric $R^1[g]$ does not vanish because the $rr$-component of the
metric contributes.

\DERIVED{OFI naturally distinguishes the ergosphere from the event horizon
through the structure of the Riesz kernel acting on the full metric tensor.}

\subsection{The Singularity at $r=0$}

At $r \to 0$, the Schwarzschild Kretschmann scalar $K \sim r^{-6}$,
which is not integrable.  The OFI operator $I^1$ acts on $K$ via
\begin{equation}
  (I^1 K)(r) \sim \int_0^\epsilon \abs{r-r'}^{-3} (r')^{-6}\,dr',
\end{equation}
which diverges: the Riesz convolution of $r^{-6}$ with $r^{-3}$ produces
a worse singularity.  \DERIVED{OFI does not resolve the black hole singularity
at $r=0$.}  This is not a failure of the framework — it reflects the genuine
non-renormalizability of point-particle gravity.

\OPEN{Full singularity resolution requires a UV completion — string theory,
LQG, or another framework.  OFI softens the Big Bang singularity (curvature
$K \sim t^{-4}$ at $\alpha=1$ becomes integrable: $\int_0^\epsilon K\,dt
\sim \epsilon^{-3} < \infty$ for $t$-integration but not pointwise) but
does not remove it.}

%%====================================================================
\section{What Is and Is Not Derived}
\label{sec:summary}
%%====================================================================

\begin{center}
\begin{tabular}{p{0.45\textwidth} p{0.45\textwidth}}
\toprule
\textbf{Derived by OFI (no free parameters)} & \textbf{Not derived — requires input or is open} \\
\midrule
$\alpha = k/2$ for every sector from kinetic symbol & Absolute values of quark and lepton masses \\
\addlinespace
Kinetic defect $= 0$ for all fields at their natural $\alpha$ & PMNS mixing angles (commute with $I^{1/2}$, not determined) \\
\addlinespace
Mass defect $= m(\bar\psi\psi - \bar\psi_a\psi_a)$ formula & Absolute neutrino masses (depend on $M_R$) \\
\addlinespace
Higgs SSB = imaginary$\to$real transition of OFI defect & Yukawa couplings $y_\nu$, $y_e$, $y_u$, $y_d$ \\
\addlinespace
Seesaw: $m_\nu \cdot M_R = m_D^2$ from defect balance & Right-handed Majorana mass $M_R$ \\
\addlinespace
Majorana defect $= \tfrac{1}{2}$ Dirac defect & Dark matter mass $m_X$ \\
\addlinespace
Graviton $\alpha = 1$ from EH kinetic degree & Dark matter halo profile exponent $\gamma$ \\
\addlinespace
Black hole horizon: $R^1[g] \propto (r-r_+)^{3/2}$ & Cosmological constant $\Lambda$ \\
\addlinespace
Ergosphere $\neq$ horizon (OFI defect structure) & Singularity resolution at $r=0$ \\
\addlinespace
Dark matter cannot collapse on cosmological timescales (SM-singlet, no EM channel) & Dark matter spin and statistics \\
\bottomrule
\end{tabular}
\end{center}

\begin{remark}[No genuine free parameters in the framework]
The OFI framework has no adjustable parameters.  The Riesz kernel is
determined by $\alpha$ and $n$.  The $\alpha$ values are determined by the
kinetic degree.  The defect formula follows algebraically.  The only inputs
are the SM Lagrangian parameters (masses, couplings) which are experimental.
\end{remark}

\begin{remark}[The remaining falsifiable prediction]
The cleanest prediction of this framework that is not already verified:
the dark matter halo density profile $\rho(r) \propto r^{-\gamma}$ with
a specific $\gamma$ determined by the OFI-Poisson system at $\alpha=1$.
If OFI predicts $\gamma = 1$ (NFW inner slope) or $\gamma = 0$ (core),
this is a falsifiable observational test.  Computing $\gamma$ is the next
concrete task.
\end{remark}

\end{document}
